\section{Controlled prime divisors and invariant section rings}
\label{V:controlled-cone-section}

Exact order preservation studies one valuation at a time.  We now give a
different finite-generation argument, in which the entire character of a
section is retained.  The geometric input is a smooth connected rationally
connected projective variety $V$, an algebraic torus $T$ acting on $V$, and
a smooth semipositive metric on $L=-K_V$.  We always use the natural
linearization of $L$.  Rational connectedness supplies
$H^i(V,\mathcal O_V)=0$ for $i>0$ and $\operatorname{Pic}^0(V)=0$.
Put $n=\dim V$.
If $V$ is a point, its constant section proves the desired nonvanishing;
assume below that $\dim V>0$.

Our first aim is to put a finite bound on a collection of divisor
classes.  Call a prime divisor $D$ on $V$ \emph{controlled} if
$aL-D$ is pseudoeffective for some integer $a>0$.  An effective integral
divisor supported on controlled primes is itself dominated by a positive
multiple of $L$: add the finitely many domination relations, with their
multiplicities.  Conversely, each component of a dominated effective
integral divisor is controlled.

We retain linearizations in this bound.  Write
$\operatorname{Pic}^T(V)_{\mathbb Q}$ for the group of linearized line
bundles tensored with $\mathbb Q$.  It is finite dimensional.  Indeed,
$\operatorname{Pic}^0(V)=0$ makes the ordinary rational Picard group
finite dimensional, and the kernel of forgetting a linearization is the
character lattice of $T$.  Every line bundle becomes linearizable after
a positive power.  To see this first for an ample line, a very ample
power gives a projective representation of the torus, which lifts after
a further power.  Write an arbitrary line as a difference of ample lines
to obtain the general assertion.

For an invariant rational divisor $D$, denote its divisor linearization
by $[D]^T$; its canonical rational section is invariant.  We regard a
character $\eta$ as the corresponding linearization of the trivial line.
With the left action on sections, an eigensection of a linearized line
$B^T$ of character $\eta$ has zero divisor of linearized class
$B^T-\eta$.

\begin{proposition}[Finite cone of controlled invariant divisors]
\label{V:controlled-cone}
Let $V$ be smooth, connected, rationally connected and projective over
$\mathbb C$.  Suppose that $L=-K_V$ has a smooth semipositive metric,
and let an algebraic torus $T$ act on $V$.  There is a rational
polyhedral cone in $\operatorname{Pic}^T(V)_{\mathbb R}$ generated by
classes of invariant effective rational divisors and containing $[D]^T$
for every invariant prime $D$ for which $aL-D$ is pseudoeffective for
some integer $a>0$.
\end{proposition}

The generators are effective but need not be controlled.  To prove
Proposition~\ref{V:controlled-cone}, we first show that all but finitely
many controlled primes come from divisors on a common base.  We then
construct an adjoint ring whose invariant generators control the
rational pullbacks of every invariant effective divisor on that base.
An auxiliary boundary makes finite generation possible and disappears
under rational pullback to $V$.

\subsection{The field of functions with controlled poles}

Let $P$ consist of zero and the rational functions whose poles are
supported on controlled primes of $V$.  Sums and products stay in $P$.
If $u\ne0$ belongs to $P$, its zero divisor is linearly equivalent to
its polar divisor, hence is dominated by a multiple of $L$.  Thus
$u^{-1}\in P$ as well.  This proves that $P$ is a $T$-stable subfield
of $\mathbb C(V)$.

The field is relatively algebraically closed.  In fact, if $u$ is
algebraic over $P$ and had a pole at an uncontrolled prime, all
coefficients of a monic equation for $u$ would have nonnegative order
there.  Its leading power would then have strictly smaller order than
every other term, which is impossible.  Also $P$ is finitely generated
over $\mathbb C$.  For completeness, choose a transcendence basis
$t_1,\ldots,t_d$ for $P/\mathbb C$ and extend it to a transcendence
basis of $\mathbb C(V)$.  The algebraic closure of
$\mathbb C(t_1,\ldots,t_d)$ inside this finitely generated field is
finite over $\mathbb C(t_1,\ldots,t_d)$, so its intermediate algebraic
extension $P$ is finite too.

Choose controlled primes $D_1,\ldots,D_s$ spanning the classes of
all controlled primes in $\operatorname{Pic}(V)_{\mathbb Q}$.  If a
controlled prime $D$ is not among them, clearing denominators and
torsion gives a principal divisor supported on $D,D_1,\ldots,D_s$
whose coefficient at $D$ is nonzero.  Its defining function belongs to
$P$.  Thus the valuation of $D$ is nontrivial on $P$, and $D$ is
vertical for any rational map with target field $P$.  In particular,
only finitely many controlled primes can be horizontal.  If
$P=\mathbb C$, this proves Proposition~\ref{V:controlled-cone} directly:
there are only finitely many controlled primes, and we take the cone of
the invariant ones.

Suppose $\operatorname{trdeg}_{\mathbb C}P>0$.
Choose a smooth projective equivariant model $Y$
of $P$ and an equivariant diagram
\begin{equation}\label{V:controlled-model}
 V\xleftarrow{\pi}U\xrightarrow{f}Y,
 \qquad \mathbb C(Y)=P,
\end{equation}
with $U$ smooth projective and $\pi$ birational.  We require that the
strict transform of each prime of $V$ either dominate $Y$ or map onto a
prime divisor of $Y$.  Here is a construction with this property.

Finitely many generators of $P$ have a common controlled polar bound
$N\ge0$.  A sufficiently large full linear system of a multiple of
$N$ contains these generators as ratios, and every ratio of its sections
still has controlled poles.  After linearizing a multiple, its image is
an equivariant projective model with function field $P$.  Over a dense
open, the closures in $V$ of the fibers form a flat family.  Resolve
equivariantly the graph of its rational map to the Hilbert scheme and
pull back the universal family.  This uses the Hilbert-family and
flattening constructions of Grothendieck and Raynaud--Gruson
\cite{Grothendieck1961Hilbert,RaynaudGruson1971}.
Flatness with fixed Hilbert polynomial
bounds fiber dimension by $\dim V-\dim Y$.  The main component is
birational to $V$.  A divisor coming from $V$ therefore cannot map to
codimension at least two: that would force its generic fiber dimension
to exceed this bound.  Resolve the total space equivariantly; strict
transforms retain the asserted images.  Relative algebraic closedness of
$P$ makes $f$ have connected fibers by Stein factorization.  Pullback of
holomorphic one-forms, followed by birational invariance, gives
$\operatorname{Pic}^0(Y)=0$.

Define rational pullback of divisors and linearized line bundles by
\[
 p^\diamond=\pi_*f^*.
\]
For line bundles, one can equivalently pull back on the big open set
where $V\dashrightarrow Y$ is defined and extend reflexively across
codimension two.  Since $V$ is smooth, this extension is a line bundle;
the equivariant structure extends uniquely.  If
$\sigma:Y'\to Y$ is a smooth birational modification, resolve the
diagram above it and use the same notation.  Then $p^\diamond$
commutes with $\sigma^*$ and kills $\sigma$-exceptional divisors,
including their divisor linearizations.  Indeed, a prime of $V$ above
such a divisor would have image of codimension at least two on $Y$,
contrary to the model property.  That property persists over $Y'$:
a horizontal prime remains horizontal, and a prime mapping onto a
divisor of $Y$ maps onto its strict transform, because $\sigma$ is
an isomorphism at its generic point.

We can now reduce the proposition to a statement about divisors on
$Y$.  On a nonempty Zariski open subset of $Y$, the morphism $f$ is
smooth with geometrically connected, hence geometrically integral,
fibers.  If the generic point of a prime $Q$ lies in this open, exactly
one prime of $U$ dominates $Q$, and its multiplicity in $f^*Q$ is one.
Every other component of $f^*Q$ maps into a smaller subset of $Y$;
by the model property it is exceptional over $V$ and pushes to zero.
Consequently, if a prime $D$ of $V$ maps onto such a $Q$, then
\[
 p^\diamond Q=D.
\]
The controlled primes not covered by this statement form a finite
set: they are the finitely many horizontal controlled primes already
identified and the finitely many primes over the excluded base divisors.
If $D$ is invariant, its image $Q$ is invariant, and the equality also
holds for their divisor linearizations.

It therefore remains to find a rational polyhedral cone of invariant
effective classes containing $p^\diamond[R]^T$ for every invariant
effective integral divisor $R$ on $Y$.  Adding the classes of the
finitely many invariant controlled exceptions will finish the proof.
We may construct this cone on any smooth equivariant modification
$\sigma:Y'\to Y$: the effective total transform $\sigma^*R$ has its
canonical divisor linearization, and
$p^\diamond[\sigma^*R]^T=p^\diamond[R]^T$.

\begin{lemma}\label{V:controlled-rank-metric-data}
For the diagram \eqref{V:controlled-model} and any very ample divisor
$H$ on $Y$, some integer $a>0$ satisfies
\begin{equation}\label{V:controlled-domination}
 a\pi^*L-f^*H\quad\text{pseudoeffective}.
\end{equation}
If $E=K_{U/V}$, then
\begin{equation}\label{V:controlled-rank}
 \operatorname{rank}f_*\mathcal O_U(E)=1.
\end{equation}
Moreover there is an invariant effective integral divisor $J$ with
\begin{equation}\label{V:controlled-J}
 (f_*\mathcal O_U(E))^{**}=\mathcal O_Y(J),
 \qquad p^\diamond[J]^T=0,
\end{equation}
where the first identification uses the canonical inclusion of constants.
\end{lemma}

\begin{proof}
The homogeneous coordinate ratios for $H$ belong to $P$, so they have
a common controlled polar bound $N$.  On $U$ the corresponding sections
have a common factor and define $f$; hence
$\pi^*N-f^*H$ is linearly equivalent to an effective divisor.
Combine this with the pseudoeffectivity of $aL-N$.

Twist $f_*\mathcal O_U(E)$ by a large multiple $bH$ so that its global
sections span its generic fiber.  Divide these sections by the Jacobian
section of $E$ and by a fixed rational frame with divisor $bH$.
The resulting rational functions on $U$ have poles bounded by
$E+f^*(bH)$.  Their pushdowns to $V$ have no poles contributed by $E$;
the remaining polar divisor is bounded by $p^\diamond(bH)$.
The preceding common-factor argument shows that this effective divisor
is controlled.  Thus all these functions belong to $P=\mathbb C(Y)$,
and their span over $P$ has dimension at most one.  The Jacobian section
is nonzero, proving \eqref{V:controlled-rank}.

The inclusion of constants makes the rank-one reflexive hull
$\mathcal O_Y(J)$ for an effective integral divisor $J$.  Every
construction is equivariant, so $J$ is invariant.  If a base prime $Q$
has a component above it that is not $\pi$-exceptional, $E$ has
coefficient zero there.  A rational function on the base with a pole at
$Q$ cannot therefore represent a local section of
$f_*\mathcal O_U(E)$.  Its coefficient in $J$ is zero.  The remaining
components of $J$ have only exceptional preimages, which proves the
second identity in \eqref{V:controlled-J}.
\end{proof}

The field construction has now supplied a rank-one adjoint image and
an exceptional correction $J$.  We next use smooth semipositivity on
$V$ to produce a strictly positive metric on $-K_Y+J$.  Its integrable
object will be the canonical section of $J$, which is the precise input
needed for weighted approximation.

\subsection{The weighted metric and its disappearing boundary}

By \eqref{V:controlled-domination}, $\pi^*L$ has a singular metric
whose curvature dominates a positive multiple of $f^*\omega_H$.
Mix its weight with the pulled-back smooth semipositive weight, using a
sufficiently small positive coefficient for the singular term.
Skoda's local integrability theorem \cite{Skoda1972} and compactness
give a metric $h$
such that
\begin{equation}\label{V:controlled-upstairs-metric}
 \Theta_h(\pi^*L)\ge\epsilon f^*\omega_H,
 \qquad \mathcal J(h)=\mathcal O_U,
 \qquad \epsilon>0.
\end{equation}

The identity
\[
 K_{U/Y}+\pi^*L=E-f^*K_Y
\]
identifies the reflexive hull of the adjoint direct image with
$-K_Y+J$.  Apply the singular direct-image metric theorem of
P\u aun--Takayama \cite[Theorem~3.3.5]{PT2018}.  Its multiplier-ideal
inclusion is an isomorphism here by
\eqref{V:controlled-upstairs-metric}.  The rank-one fiber integral
therefore gives a singular metric with weights $\varphi$ on
$-K_Y+J$, extending across divisors and then across codimension two.
It has curvature at least $\epsilon\omega_H$: locally subtract a
potential of this form from the coefficient weight before applying
positivity, and restore it afterward.

For clarity, its required integrability follows directly from the same
integral.  The section induced by the Jacobian section of $E$ is
represented by $s_J$ in a local frame of the reflexive hull.  Integrating
its squared fiber norm against base coordinate volume gives the integral
upstairs of the corresponding canonical-valued form with coefficient
metric $h$.  This is finite by $\mathcal J(h)=\mathcal O_U$ and
properness.  Fiber integration on the smooth locus, which has full
measure, proves
\begin{equation}\label{V:controlled-weighted-integral}
 |s_J|^2e^{-\varphi}\in L^1_{\mathrm{loc}}(Y),
 \qquad \mathcal J(\varphi)\supseteq\mathcal O_Y(-J).
\end{equation}
This is the weighted input of \cite[Theorem~\textup{(Weighted anticanonical approximation)}]{CompanionE}.  It does not require an unweighted
integrability assertion for $e^{-\varphi}$.

We record the complete output of that theorem at this use.  A power of
an ample line on $Y$ is $T$-linearizable, since
$\operatorname{Pic}^0(Y)=0$; the zero divisor of an eigensection of a
very ample power is invariant and ample.  Choose a small positive
rational multiple $A$ of this divisor so that $B=-K_Y+J-A$ still has
a strict curvature bound.  Resolve the full
system of a sufficiently divisible $rB$ together with $J$ by an
equivariant sequence of smooth blowups $\sigma:Y'\to Y$ and write
\[
 \sigma^*|rB|=G+|\mathrm{Mov}|,
\]
with the mobile system base-point free.  Weighted approximation gives
\begin{equation}\label{V:controlled-signed-boundary}
 \Psi=G/r-\sigma^*J-K_{Y'/Y},\qquad
 \operatorname{coeff}\Psi<1.
\end{equation}
Writing $\Psi=\Psi^+-\Psi^-$ with disjoint positive and negative
parts, put $D_0=\Psi^-$.  Then
\begin{equation}\label{V:controlled-adjoint}
 D_0\sim_{\mathbb Q}K_{Y'}+\Psi^+
             +\mathrm{Mov}/r+\sigma^*A.
\end{equation}
Here $\Psi^+$ is an effective simple-normal-crossings klt boundary,
and the last two summands are nef and big together.  The support of
$D_0$ is contained in that of $\sigma^*J$ and the exceptional divisor
of $\sigma$.  The divisors $G,K_{Y'/Y},\Psi^+,D_0$ are invariant and
the free system $|\mathrm{Mov}|$ is torus-stable.  Identity
\eqref{V:controlled-adjoint} is an ordinary rational linear equivalence;
we construct the needed linearized classes separately below.

The support statement has the crucial additional consequence
\begin{equation}\label{V:controlled-D0-zero}
 p^\diamond[D_0]^T=0.
\end{equation}
Indeed $p^\diamond$ kills the effective divisor $\sigma^*J$ by
\eqref{V:controlled-J}, and it kills every $\sigma$-exceptional prime
by the model property.  An effective divisor supported on this union
has zero pushdown.  Its canonical invariant section also pushes to the
constant rational section, so the equality includes the linearization.
No character shift is lost in \eqref{V:controlled-D0-zero}.

\subsection{The multigraded ring and the finite cone}

We now build the cone of base pullbacks required above.  Let $R_0$
be any invariant effective integral divisor on $Y'$.  Since $[D_0]^T$
disappears under $p^\diamond$, the divisor $D_0+tR_0$ for a small
rational $t>0$ has the same rational pullback class as $tR_0$.
We will place these divisors in a finite multigraded ring and express
their invariant sections through its generators.  This is why we need
a whole neighborhood of $[D_0]^T$ in the linearized Picard space,
rather than just the single adjoint ring of $D_0$.

Because $\sigma$ is a sequence of smooth blowups, there is an effective
exceptional rational divisor $F$ with $-F$ relatively ample.  One
constructs it successively by giving sufficiently small positive
coefficients to the later exceptional divisors.  For small rational
$\delta>0$, $\sigma^*A-\delta F$ is ample and
$\Psi^++\delta F$ is still a klt boundary with simple normal
crossings.  If $\sigma$ is the identity, take $F=0$.  Thus
\eqref{V:controlled-adjoint} expresses $D_0$ as $K_{Y'}$ plus that
boundary and an ample rational class.

Choose a fixed small ample rational part $A_*$.  Every rational class
in a sufficiently small neighborhood of the underlying class of $D_0$
can be written
\[
 K_{Y'}+A_*+\Delta_i,
\]
with $\Delta_i\ge0$ and klt.  For a finite set of classes, choose
general members of sufficiently divisible very ample systems for the
remaining ample parts, divided by their multiples.  They can be chosen
transverse to one another and to the fixed boundary, with coefficients
less than one.  Consequently their union has simple normal crossings.
Small character changes only change linearizations, so the same
construction gives a neighborhood in
$\operatorname{Pic}^T(Y')_{\mathbb R}$.

Choose rational vertices $U_1^T,\ldots,U_s^T$ of a polytope with
$[D_0]^T$ in its interior.  After a common multiple $k$, let
$\mathcal E_i$ be integral linearized representatives of $kU_i^T$.
The ordinary ring
\begin{equation}\label{V:controlled-ring}
 R=\bigoplus_{\boldsymbol\ell\in\mathbb N^s}
 H^0\!\left(Y',\bigotimes_{i=1}^s
                      \mathcal E_i^{\otimes\ell_i}\right)
\end{equation}
is finitely generated by the multigraded adjoint-ring theorem of
Cascini--Lazi\'c \cite[Theorem~A]{CasciniLazic2012}; compare
\cite[Theorem~3.2(1)]{CortiLazic2013}.  We take $k$ sufficiently
divisible to turn the rational linear equivalences above into
isomorphisms of line bundles.  This identifies \eqref{V:controlled-ring}
with a divisible subring of the stated adjoint ring.  The linearizations
give an algebraic $T$-action on $R$, independently of the general
boundary representatives used to prove finite generation.

Here is the invariant-ring step explicitly.  Enlarge finitely many
multihomogeneous generators of $R$ to their finitely many character
components.  They define an equivariant multigraded surjection from a
polynomial ring to $R$.  Projection to weight zero is exact for a torus,
so the invariant polynomial ring surjects onto $R^T$.  If the variable
weights are $\eta_1,\ldots,\eta_N$, its invariant monomials have
exponents
\[
 (a_1,\ldots,a_N)\in\mathbb N^N,
 \qquad \sum a_i\eta_i=0.
\]
These are the lattice points of a rational polyhedral cone; Gordan's
lemma makes their semigroup finitely generated.  Thus $R^T$ has
finitely many nonzero multihomogeneous generators $u_1,\ldots,u_b$.

Let $\mathcal C_0$ be the cone generated by the rational pullbacks
$p^\diamond[\operatorname{div}(u_j)]^T$.  These are classes of
invariant effective divisors on $V$.  We claim that $\mathcal C_0$
contains $p^\diamond[R_0]^T$ for \emph{every} invariant effective
integral divisor $R_0$ on $Y'$.  Choose a small rational $t>0$ such
that $[D_0]^T+t[R_0]^T$ belongs to the vertex polytope.  Express it as
a rational convex combination of the $U_i^T$.  Clearing denominators
and torsion gives an equivariant isomorphism between a positive
multiple of $D_0+tR_0$ and one multidegree in
\eqref{V:controlled-ring}.  Its canonical section gives a nonzero
element of $R^T$ in that multidegree.  Express it as a polynomial in
the $u_j$ and retain the terms of that multidegree.  At least one
nonzero monomial occurs, and its degree is a nonnegative integral sum
of generator degrees.  Thus the corresponding linearized class lies
in their cone.  Only equality of homogeneous degrees is used here;
the selected monomial need not have the same zero divisor as the
original section.  Apply $p^\diamond$, use
\eqref{V:controlled-D0-zero}, and divide by the positive factor
containing $t$.  This proves the claim.

Apply this claim to the total transforms of invariant effective
divisors on the original base $Y$.  Compatibility of $p^\diamond$
with $\sigma^*$ gives the required cone of their rational pullbacks.
Add the invariant effective classes of the finitely many controlled
primes excluded in the initial comparison with base divisors.  The
resulting rational polyhedral cone proves
Proposition~\ref{V:controlled-cone}.

\subsection{A weight bound and character separation}

The cone gives finitely many effective generators.  To obtain an
invariant anticanonical section, we must still put the \emph{specified
natural linearization} of $L$ in their cone.  The following fixed-point
calculation supplies the sign needed for separation.  It is a bound on
some cohomology weight, rather than the weight-zero index theorem.

\begin{lemma}\label{V:bounded-above-weight}
Let $V$ be smooth connected projective with
$H^i(V,\mathcal O_V)=0$ for $i>0$, and let
$v:\mathbb G_m\to\operatorname{Aut}(V)$ be an algebraic action.
Give $L=-K_V$ its natural linearization.  There is a constant $C$ such
that, for every sufficiently large integer $m$, some
$H^q(V,mL)$ has a $v$-weight at most $C$.
\end{lemma}

\begin{proof}
If the action is trivial, all weights are zero and the ordinary
Riemann--Roch polynomial $\chi(V,mL)$ has constant term one, so the
assertion follows.  Otherwise use the holomorphic fixed-point formula
of Atiyah--Bott, in its fixed-component equivariant index form
\cite{AtiyahBott1968,AtiyahSegal1968,AtiyahSinger1968}.
For a fixed component $Z$, let $a$ be the nonzero
weights of its normal tangent bundle, with Chern roots $x_{a,j}$.
The anticanonical fiber weight is $b_Z=\sum_{a,j}a$.  Its contribution
to the alternating cohomology character is
\[
 t^{mb_Z}\int_Z e^{m c_1(L|_Z)}\operatorname{td}(Z)
          \prod_{a,j}(1-t^{-a}e^{-x_{a,j}})^{-1}.
\]
The integral is a finite cohomological expansion.  Denote the alternating
sum of the cohomology characters by $\chi_v(V,mL)$.  Grouping by $b_Z$
writes it as
\begin{equation}\label{V:weight-windows}
 \chi_v(V,mL)=\sum_b t^{mb}R_b(m,t),
\end{equation}
where each $R_b$ is polynomial in $m$ with rational functions of $t$
as coefficients.

At $m=0$, the left side is one.  As $t\to0$, each component having
a positive normal tangent weight has contribution tending to zero;
the other components have $b_Z\le0$.  Hence at least one $b\le0$
has $R_b(0,t)\ne0$.  Multiply \eqref{V:weight-windows} by one fixed
nonzero Laurent polynomial clearing the denominators of every
coefficient.  For the selected $b$, the exponents lie in
$mb+[-C_0,C_0]$.  Some coefficient is a nonzero polynomial in $m$,
so this block is nonzero for all large $m$.  The finitely many windows
for different $b$ are disjoint for large $m$.  Thus the multiplied
character has a nonzero exponent bounded above independently of $m$.
Removing a fixed Laurent-polynomial factor changes the minimum
exponent only by a fixed number.  The original virtual character
therefore has such an exponent too, and it occurs in at least one
actual cohomology group.
\end{proof}

We now obtain an independent proof of the rationally connected
special case of Theorem~\ref{V:invariant-main}. Its additional output
is the controlled-divisor cone just constructed. The proof uses that
cone and character separation, and will supply the fiberwise section
in the locally constant MRC application of Section~\ref{V:sec-mrc-descent}.

\begin{theorem}\label{V:controlled-invariant-nonvanishing}
Let $V$ be a smooth connected rationally connected projective complex
variety with smoothly semipositive $-K_V$.  For every algebraic torus
$T$ acting on $V$, there is an integer $m>0$ such that
\[
 H^0(V,-mK_V)^T\ne0
\]
with the natural anticanonical linearization.
\end{theorem}

\begin{proof}
Let $\mathcal C$ be the cone in Proposition~\ref{V:controlled-cone}.
Suppose $L^T\notin\mathcal C$.  Rational polyhedral separation gives
a rational functional $\lambda$ on the linearized Picard space with
\[
 \lambda(L^T)>0,\qquad \lambda(\mathcal C)\le0.
\]
After positive rescaling, its restriction to characters is pairing with
a cocharacter $v:\mathbb G_m\to T$, possibly zero.
Lemma~\ref{V:bounded-above-weight} gives cohomology classes with
$v$-weights uniformly bounded above and with unbounded positive $m$.

Average a K\"ahler form and the smooth metric over the maximal compact
torus.  Smooth-coefficient hard Lefschetz then gives the equivariant
surjection
\[
 H^0(V,\Omega_V^{n-q}\otimes(m+1)L)
       \longrightarrow H^q(V,K_V+(m+1)L)=H^q(V,mL).
\]
This is the smooth case of the Demailly--Peternell--Schneider theorem
\cite[Theorem~0.1]{DPS2001}.
The coefficient shift $m+1$ and its equivariant identification use the
natural action on $L$.  Pass to a subsequence with fixed
$p=n-q$, and choose $T$-eigensections $\alpha_m$ of
$\Omega_V^p\otimes(m+1)L$ whose $v$-weights are bounded above.

Take the generic span of their untwisted lines in $\Omega_V^p$,
of rank $s$, and saturate its determinant to a line bundle
$M\subseteq\bigwedge^s\Omega_V^p$.  This line inherits a
$T$-linearization.  The inverse-determinant pseudoeffectivity theorem
of Lazi\'c--Matsumura--Peternell--Tsakanikas--Xie
\cite[Theorem~4.1]{LMPTX2023}, following Ou's generic nefness theorem
\cite[Theorem~1.4]{Ou2023}, gives $-M$ pseudoeffective.  Its hypotheses apply because $V$ is
smooth projective and $-K_V$ is nef; the exterior construction is a
direct summand of a cotangent tensor power, and saturation of the line
persists in that summand.  If $p=0$, $M=\mathcal O_V$ instead.

Choose finitely many $\alpha_m$ forming a generic basis.  Each further
nonzero $\alpha_m$ can be completed to a generic basis using this
fixed set.  Their wedge gives a nonzero eigensection of $M+kL$, for
unbounded positive $k$, whose $v$-weight is still bounded above.  The
wedge lands in $M+kL$ by saturation in codimension one and reflexive
extension.  Let $D_k$ be its divisor and $\eta_k$ its character.
Then $kL-D_k$ has class $-M$, so each component of $D_k$ is controlled.
Each is also invariant, because $T$ is connected and $D_k$ is the
zero divisor of an eigensection.  Consequently $[D_k]^T\in\mathcal C$.
On the other hand,
\[
 \lambda([D_k]^T)=\lambda(M^T)+k\lambda(L^T)
                              -\langle\eta_k,v\rangle>0
\]
for large $k$, a contradiction.

Thus $L^T\in\mathcal C$.  A real nonnegative expression in its
finitely many rational generators admits a rational nonnegative
expression.  Clear denominators and torsion in the linearized Picard
group and multiply the canonical invariant sections of these effective
divisors.  The product is nonzero since $V$ is integral, and is a
section of a positive multiple of the naturally linearized $L$.
If $L^T$ is zero in the rational class space, clearing its torsion gives
the constant invariant section directly.
\end{proof}
